\section*{Capítulo \ref{ch:g12:equilibrium} --- O equilíbrio químico: quociente e constante}

\begin{solution}{exo:g12:equilibrium:1}
(a) $Q = \dfrac{[\ce{NH3}]^2 (c^\circ)^2}{[\ce{N2}][\ce{H2}]^3}$;
(b) $Q = \dfrac{[\ce{CH3COO-}][\ce{H3O+}]}{[\ce{CH3COOH}]\,c^\circ}$ (a
água, o \omterm{def:g6:solutions-and-solubility:solute}{solvente}, fica de fora);
(c) $Q = [\ce{Ca^{2+}}][\ce{CO3^{2-}}]/(c^\circ)^2$ (o sólido fica de fora);
(d) $Q = \dfrac{[\ce{FeSCN^{2+}}]\,c^\circ}{[\ce{Fe^{3+}}][\ce{SCN-}]}$.
\end{solution}

\begin{solution}{exo:g12:equilibrium:2}
$\tau = 0.020/0.050 = 0.40$: não é total.
\end{solution}

\begin{solution}{exo:g12:equilibrium:3}
$Q = 2 < K$: sentido direto. $Q = 50 > K$: sentido inverso. $Q = K$: nada muda.
\end{solution}

\begin{solution}{exo:g12:equilibrium:4}
A composição é constante, mas as reações direta e inversa continuam na
mesma velocidade: as \omterm{def:g7:atoms-and-molecules:molecule}{moléculas} não param de reagir nos dois sentidos.
\end{solution}

\begin{solution}{exo:g12:equilibrium:5}
Não, $K$ não muda (o \omterm{def:g10:reaction-progress-table:system}{estado final} é o mesmo); o tempo para chegar ao
equilíbrio fica mais curto.
\end{solution}

\begin{solution}{exo:g12:equilibrium:6}
$\dfrac{x^2}{(2-x)^2} = 4$, logo $\dfrac{x}{2-x} = 2$ e
$x = \qty{1.33}{mol}$: \qty{0.67}{mol} de ácido e de \omterm{def:g11:functional-groups:alcohol}{álcool},
\qty{1.33}{mol} de \omterm{def:g11:functional-groups:carboxylic-acid}{éster} e de água (de novo dois terços).
\end{solution}

\begin{solution}{exo:g12:equilibrium:7}
$Q = (0.5 \times 0.5)/(0.5 \times 0.5) = 1 < 4$: sentido direto, forma-se
mais \omterm{def:g11:functional-groups:carboxylic-acid}{éster}.
\end{solution}

\begin{solution}{exo:g12:equilibrium:8}
Cerca de \qty{0.52}{mol} do lado do ácido e \qty{0.79}{mol} do lado do
\omterm{def:g11:functional-groups:carboxylic-acid}{éster}. Depois de cerca de \qty{100}{h}, as duas curvas estão sobre a
linha tracejada: equilíbrio.
\end{solution}

\begin{solution}{exo:g12:equilibrium:9}
Com $y$ \omterm{def:g10:the-mole:amount}{mol} hidrolisados: $\dfrac{y^2}{(1-y)^2} = \dfrac{1}{4}$, logo
$\dfrac{y}{1-y} = \dfrac{1}{2}$, $y = \frac{1}{3}$: $\frac{2}{3}$ \omterm{def:g10:the-mole:amount}{mol} de
\omterm{def:g11:functional-groups:carboxylic-acid}{éster} e de água, $\frac{1}{3}$ \omterm{def:g10:the-mole:amount}{mol} de ácido e de \omterm{def:g11:functional-groups:alcohol}{álcool}. A mesma \omterm{def:g6:pure-substances-and-mixtures:mixture}{mistura}
final que a da esterificação.
\end{solution}

\begin{solution}{exo:g12:equilibrium:10}
Um sólido é uma fase pura, cuja ``concentração'' não muda: ele fica de
fora de $Q$. Acrescentar mais sólido não muda $Q$, logo não desloca o
equilíbrio (enquanto houver sólido presente).
\end{solution}

\begin{solution}{exo:g12:equilibrium:11}
$K$ fixa a razão entre o dióxido de carbono no gás e o dióxido de
carbono na água. Com a garrafa aberta, o gás acima da água escapa para o
\omterm{def:g4:air-a-mixture-of-gases:air}{ar}: o dióxido de carbono no gás diminui, $Q < K$, e mais gás sai da
água, até quase não sobrar nenhum.
\end{solution}

\begin{solution}{exo:g12:equilibrium:12}
No equilíbrio: ácido e \omterm{def:g11:functional-groups:alcohol}{álcool} $\frac{1}{3}$, \omterm{def:g11:functional-groups:carboxylic-acid}{éster} e água
$\frac{2}{3}$. Retirar a metade da água deixa $\frac{1}{3}$:
$Q = \dfrac{(2/3)(1/3)}{(1/3)(1/3)} = 2 < 4$, sentido direto. Com mais $y$
reagindo: $(2/3 + y)(1/3 + y) = 4(1/3 - y)^2$, isto é,
$27y^2 - 33y + 2 = 0$, $y = 0.064$: o \omterm{def:g11:functional-groups:carboxylic-acid}{éster} sobe para \qty{0.73}{mol}.
\end{solution}

\begin{solution}{exo:g12:equilibrium:13}
O quociente da reação inversa tem o numerador e o denominador trocados:
$Q' = 1/Q$, logo, no equilíbrio, $K' = 1/K$. Hidrólise:
$1/4 = 0.25$.
\end{solution}

\begin{solution}{exo:g12:equilibrium:14}
$Q = (2 \times 2)/(1 \times 1) = 4 = K$: a \omterm{def:g6:pure-substances-and-mixtures:mixture}{mistura} já está em
equilíbrio; a sua composição não muda.
\end{solution}

\begin{solution}{exo:g12:equilibrium:15}
$n - 0.95 = 0.9025 / (4 \times 0.05) = 4.51$, logo $n = \qty{5.5}{mol}$
de \omterm{def:g11:functional-groups:alcohol}{álcool} por \omterm{def:g10:the-mole:amount}{mol} de ácido.
\end{solution}

\begin{solution}{pb:g12:equilibrium:1}
\textbf{1.} \ce{CH3-COOH + CH3-CH2-OH <=> CH3-COO-CH2-CH3 + H2O}.

\textbf{2.} Ela não é total: para num equilíbrio em que os dois
sentidos continuam.

\textbf{3.} $Q = \dfrac{n_{\text{éster}}\, n_{\text{água}}}{n_{\text{ácido}}\, n_{\text{álcool}}}$.

\textbf{4.} $x_{\max} = \qty{1}{mol}$; $x_f = \frac{2}{3}$ \omterm{def:g10:the-mole:amount}{mol}.

\textbf{5.} $\tau = 0.67$.

\textbf{6.} Ácido e \omterm{def:g11:functional-groups:alcohol}{álcool}, $\frac{1}{3}$ \omterm{def:g10:the-mole:amount}{mol} cada; \omterm{def:g11:functional-groups:carboxylic-acid}{éster} e água,
$\frac{2}{3}$ \omterm{def:g10:the-mole:amount}{mol} cada.

\textbf{7.} $K = \dfrac{(2/3)^2}{(1/3)^2} = 4$.

\textbf{8.} Um terço hidrolisado: \omterm{def:g11:functional-groups:carboxylic-acid}{éster} e água $\frac{2}{3}$, ácido e
\omterm{def:g11:functional-groups:alcohol}{álcool} $\frac{1}{3}$: a mesma \omterm{def:g6:pure-substances-and-mixtures:mixture}{mistura}, $Q_{eq} = 4$.

\textbf{9.} É a mesma reação à mesma temperatura, e $K$ só depende da
temperatura.

\textbf{10.} Ácido $1 - x$; \omterm{def:g11:functional-groups:alcohol}{álcool} $3 - x$; \omterm{def:g11:functional-groups:carboxylic-acid}{éster} $x$; água $x$.

\textbf{11.} $x^2 = 4(1 - x)(3 - x) = 4(3 - 4x + x^2)$, logo
$3x^2 - 16x + 12 = 0$.

\textbf{12.} $x = (16 \pm \sqrt{256 - 144})/6 = (16 \pm 10.58)/6$: 0.903
ou 4.43; só $x_f = \qty{0.903}{mol}$ fica entre 0 e 1.

\textbf{13.} \qty{90}{\percent}.

\textbf{14.} $M = \qty{88.0}{g/mol}$: $0.903 \times 88.0 = \qty{79.5}{g}$.

\textbf{15.} Cada retirada de água diminui o numerador de $Q$, que fica
abaixo de $K$: o sistema volta a avançar.

\textbf{16.} A reação continuaria até o ácido acabar:
\qty{100}{\percent}.

\textbf{17.} O etanol é mais barato, e o seu excesso é facilmente
separado do \omterm{def:g11:functional-groups:carboxylic-acid}{éster} (e reciclado).

\textbf{18.} Só o tempo: uma reação que quase não libera calor tem uma
constante que quase não varia com a temperatura, e aquecer só a faz
chegar mais cedo ao equilíbrio.

\textbf{19.} Cerca de \qty{90}{\percent} do ácido.
\end{solution}
